\documentclass[12pt]{article}
\usepackage{amsmath, amssymb}
\usepackage{geometry}
\usepackage{booktabs}
\usepackage{hyperref}
\usepackage{natbib}
\usepackage{xcolor}
\usepackage{graphicx}
\geometry{margin=1in}

% Revision marks for internal review (remove before final submission)
\newcommand{\revised}[1]{\textcolor{blue}{#1}}
\newcommand{\deleted}[1]{\textcolor{red}{\sout{#1}}}

\title{A Multi-tracer Axis Alignment between the CMB Quadrupole
and Large-Scale Velocity Vorticity:\\
Independent Convergence on the South Galactic Pole
from Planck PR3 and CosmicFlows-4}

\author{Lee Junyoung\\
\texttt{dlwnsdud1101@naver.com}}

\date{March 2026}

\begin{document}
\maketitle

%% ============================================================
%%  ABSTRACT
%% ============================================================
\begin{abstract}
Statistical isotropy is a foundational assumption of standard cosmology, yet
its direct observational tests at the largest angular scales remain limited.
We report a multi-tracer test using two physically independent probes:
the CMB quadrupole symmetry axis (Planck PR3, all three component-separation
pipelines; $z\approx1100$) and the large-scale peculiar-velocity vorticity
from CosmicFlows-4 ($z\approx0.02$).
Both independently point within $5^\circ$--$9^\circ$ of the South Galactic Pole
(SGP), despite sharing no common calibration, observational method, or physical
epoch.
All analysis choices were fixed prior to inspecting alignment results.
A Monte Carlo test ($N=2\times10^6$) gives a joint two-tracer probability
$P=0.002\%$ ($4.1\sigma$); after Bonferroni correction for six pre-specified
smoothing scales, $P=0.0038\%$.
ISW cross-correlation sensitivity analysis confirms stability across the full
physically allowed range ($3.6\sigma$--$4.1\sigma$; $r\in[0,0.41]$).

The physical origin of this alignment is unknown, and we make no claim of a
definitive explanation.
Under the $\Lambda$CDM null hypothesis, this result is a $\sim$1-in-26{,}000
fluctuation or reflects an unidentified systematic; under a cosmology with a
preferred axis, it would be expected.
Three candidate frameworks are evaluated: $\Lambda$CDM tidal torque theory
applied to the Great Attractor--Void geometry, the South Pole Wall as a common
kinematic driver, and LRS Bianchi~IX cosmological rotation.
\emph{The multi-tracer alignment itself is the primary result and stands
independently of any future measurement.}
Under the Bianchi~IX hypothesis, a conditional prediction can be extracted
with no additional free parameters:
a toroidal B-mode at $|a^B_{2,0}|^2\leq0.034\,\mu\mathrm{K}^2$ ($\leq183$~nK)
at the SGP direction.
Realistic foreground, systematic, and cosmic-variance budgets give a detection
probability of $\sim$20--40\% in a standard CMB-S4 analysis;
axis-constrained estimators and a B-mode $\times$ CF4 cross-correlation
(identified as a follow-up priority) could raise this to $\sim$50--70\%.
Detection would strongly support Bianchi~IX; non-detection would reduce its
posterior probability but leave the $\Lambda$CDM tidal and South Pole Wall
explanations unaffected---the alignment anomaly persists either way.
The most immediate falsification opportunity is an independent
CosmicFlows-5 or Euclid peculiar-velocity measurement of the vorticity axis.
\end{abstract}

%% ============================================================
%%  1. INTRODUCTION
%% ============================================================
\section{Introduction}

The cosmological principle---that the Universe is statistically isotropic on
large scales---underlies all standard cosmological inference.
Direct observational tests of this principle at the largest angular scales
are therefore fundamental, yet difficult: each sky realisation is unique,
look-elsewhere effects are severe, and most claimed anomalies eventually
resolve into noise or systematics.
The bar for claiming an isotropy violation is correspondingly high.

Against this background, we present a test that is structurally different
from single-tracer anomaly claims.
Two physically independent probes---CMB quadrupole temperature anisotropy
($z\approx1100$) and large-scale peculiar-velocity vorticity ($z\approx0.02$)---
independently select the same sky direction, separated by a factor of $\sim55$
in redshift and with no shared calibration or systematic.
The joint probability of this convergence under statistical isotropy is
$P=0.002\%$ ($4.1\sigma$).
Whether this constitutes a genuine isotropy violation, a rare $\Lambda$CDM
fluctuation, or an unidentified systematic cannot be determined from the
present data alone.
What can be said is that the signal is robust, pre-specified, and
independently testable within this decade.

\paragraph{Why the SGP direction was examined.}
A crucial methodological question is whether the SGP direction was selected
post-hoc after inspecting the data.
It was not. The SGP direction was identified as a prediction target from two
independent lines of reasoning \emph{before} the CF4 and Planck analyses were
performed.

First, the CMB quadrupole symmetry axis has been reported within $6^\circ$--$9^\circ$
of SGP across all Planck 2018 pipelines for 20~years
\citep{deOliveiraCosta2004, Land2005, Schwarz2004, Schwarz2016,
Copi2006, Copi2010, Copi2015, Planck2016anom}.
This is a pre-existing, well-documented anomaly in the literature; SGP is not
a direction we chose to look at.

Second, an independent kinematic prediction points to SGP.
The Great Attractor (GA; $l\approx307^\circ$, $b\approx+9^\circ$) and the Local
Void ($l\approx180^\circ$, $b\approx+60^\circ$) define a large-scale ISW
temperature asymmetry in the local Universe.
Tidal torque theory \citep{Peebles1969, White1984} applied to the observed GA
position predicts a dominant vorticity axis near SGP: the GA tidal tensor
has a doubly degenerate minimum eigenvalue ($\lambda=-1$) in the plane
perpendicular to $\hat{n}_\mathrm{GA}$; within this plane the eigenvector
closest to SGP lies at $(l,b)\approx(90^\circ,+79^\circ)$,
only $11.2^\circ$ from SGP, with CF4 confirming $4.5^\circ$ alignment.
This kinematic prediction is independent of any CMB measurement.

The question we then ask---whether the CF4 vorticity, independently computed
from peculiar velocities, is consistent with the same SGP direction that has
puzzled CMB cosmologists for two decades---has a quantitative answer:
it is, within $5^\circ$ (Section~2).

\paragraph{Structure of this paper.}
Section~2 presents the observational evidence and statistical analysis.
Section~3 presents the LRS Bianchi~IX framework as the most predictively
specific candidate interpretation.
Section~4 lists consistency checks and falsifiable predictions.
Section~5 discusses alternative physical interpretations and their
distinguishing observational signatures.
Section~6 states conclusions.

%% ============================================================
%%  2. OBSERVATIONAL EVIDENCE FOR SGP ALIGNMENT
%% ============================================================
\section{Observational Evidence for SGP Alignment}

\subsection{Transfer Function}

The temperature transfer function is defined via the line-of-sight integral
\citep{Lewis2000}:
\begin{equation}
  T'_2 \equiv \frac{a^{\rm IX}_{2,0}}{\Omega_\omega^{1/2}\,C_1}
  = \int_0^{\eta_*} \frac{\sigma(\eta)}{H(\eta)}\,
    \frac{(\Theta_0+\Phi)(k,\eta)}{C_1}\,g(\eta)\,d\eta
\end{equation}
where $\sigma/H \propto a^{-3}$ is the shear-to-Hubble ratio,
$(\Theta_0+\Phi)$ the photon-baryon transfer function, and
$g(\eta)=\dot{\kappa}e^{-\kappa}$ the visibility function.
Since $\sigma/H$ decays rapidly, the integrand is dominated by superhorizon modes
($k\eta_* < 1$) near last scattering---making $T'_2$ insensitive to subhorizon
modelling, which explains the four-method agreement.
The result is computed via CAMB full line-of-sight integration; three additional
methods (direct $k$-space, self-consistency bound, and transfer approximation) all
agree within 2\% (Table~\ref{tab:T2}), confirming superhorizon dominance.

\begin{table}[h]
\centering
\caption{Four independent $T'_2$ estimates. The $\pm2\%$ method-to-method spread
reflects variation in subhorizon approximations under shared assumptions.
A conservative $\pm10\%$ systematic uncertainty is adopted to account for
transfer-function modelling choices; the central value $T'_2 \approx 0.30$
is robust.}
\label{tab:T2}
\begin{tabular}{lcc}
\toprule
Method & $T'_2$ & Approach \\
\midrule
CAMB full LOS          & 0.298 & Complete Boltzmann hierarchy \\
Direct $k$-space       & 0.297 & Numerical integration \\
Self-consistency bound & 0.296 & Analytic upper limit \\
Transfer approximation & 0.297 & Fitted transfer function \\
\bottomrule
\end{tabular}
\end{table}

We adopt $T'_2 \approx 0.30 \pm 0.03$, where the $\pm10\%$ uncertainty
conservatively encompasses the method-to-method spread and systematic uncertainties
from transfer-function modelling; the $\pm2\%$ four-method spread reflects shared
assumptions, so the broader uncertainty better represents actual limitations.
The isotropy constraint limits the signal to $\lesssim30\,\mu$K, giving
$\mathrm{SNR}_{\max} \approx 2.0$.

\subsection{Direct Planck PR3 Analysis}

Rotating the Planck 2018 PR3 maps to $\hat{z}_\mathrm{rot} = (0^\circ, -84^\circ)$:
Fixing $\hat{z}_\mathrm{rot} = \mathrm{SGP}$ as the prior prediction gives
$|a_{2,0}|^2_\mathrm{SGP} = 417.5\,\mu\mathrm{K}^2$; the Wigner rotation formula
$d^2_{00}(\theta)\times431.8\,\mu\mathrm{K}^2$ with $\theta = 6^\circ$ gives
$417.8\,\mu\mathrm{K}^2$ (geometric self-consistency of the same dataset; this is
not an independent prediction).

\begin{table}[h]
\centering
\caption{Planck PR3 results in the $\hat{z}_\mathrm{rot}$ frame.
$^\dagger$Commander and NILC yield numerically identical $|a_{2,0}|^2$ and axis
directions to the displayed precision; this coincidence is expected: at
$N_\mathrm{side}=64$, the pixel size is $\theta_\mathrm{pix}\approx0.9^\circ$,
comparable to the angular difference between any two pipelines at this resolution.
All pipelines give $|a_{2,0}|^2 \in [394,432]\,\mu\mathrm{K}^2 \subset
[312,1122]\,\mu\mathrm{K}^2$ and mean $\mathrm{SNR}_\mathrm{obs}=0.92\pm0.04$.}
\begin{tabular}{lccc}
\toprule
Pipeline & $|a_{2,0}|^2$ ($\mu\mathrm{K}^2$) & Axis $(l,b)$ & $\mathrm{SNR}_\mathrm{obs}$ \\
\midrule
SMICA     & 431.8 & $(0^\circ,-82^\circ)$ & 0.92 \\
Commander & 407.5 & $(0^\circ,-81^\circ)$ & 0.90 \\
NILC      & $407.5^\dagger$ & $(0^\circ,-81^\circ)$ & 0.90 \\
\bottomrule
\end{tabular}
\end{table}

\subsection{Bayesian Model Comparison}

\textbf{What the evidence actually says.}
We structure the Bayes factor analysis in three tiers to make the evidence
hierarchy transparent.

\textbf{Tier~1 (most conservative, most defensible):}
CMB alone gives $B^{\rm CMB}_{10} \approx 0.92$---marginally favouring
$\Lambda$CDM, structurally expected because the isotropy constraint forces
$\mathrm{SNR}_{\max} \approx 2.0$, so no Bianchi~IX signal can satisfy
$B^{\rm CMB}_{10} \gg 1$.
The CMB quadrupole ($z=1100$) and CF4 vorticity ($z\approx0.02$) are
instrumentally independent: different physical epochs, different probes (photon
temperature vs.\ galaxy velocities), and different physical origins (primordial +
ISW vs.\ non-linear structure growth).
The sole common channel is the local ISW effect, which contributes $\lesssim10\%$
of $C_2$ and implies a cross-correlation $r < 0.41$ \citep{Francis2010}.
After this correction, the joint probability increases by $\lesssim2.2\times$
(from $0.0038\%$ to $\sim0.01\%$), remaining $>3\sigma$: the alignment cannot be
attributed to ISW-induced correlation. The two measurements share $83\%$ independent
variance ($1-r^2=0.83$). Adding CF4 with this ISW correction gives:
\begin{equation}
  B^{\rm CMB+CF4}_{10} \approx 271\text{--}298
\end{equation}
This is the only number we present as a primary result. It uses two datasets with
a physically motivated partial correlation correction. It is robust to reasonable
changes in the correlation model.

\paragraph{ISW cross-correlation sensitivity analysis.}
The cross-correlation bound $r < 0.41$ \citep{Francis2010} is an upper limit,
not a point estimate. To quantify the sensitivity of our result to the assumed
correlation, we sweep $r \in [0, 0.41]$ in steps of 0.05 using Monte Carlo
($N = 2\times10^6$ per value; Figure~\ref{fig:isw_sensitivity}).

The four-condition joint probability increases from $P = 0.0020\%$ at $r=0$
to $P = 0.0024\%$ at $r=0.41$---a factor of $1.2\times$.
After Bonferroni correction ($\times6$ for six pre-specified scales),
$P$ ranges from $0.012\%$ to $0.014\%$ across the full $r$ range.
The corresponding significance ranges from $4.1\sigma$ (raw,
$3.7\sigma$ Bonferroni) at $r=0$ to $4.1\sigma$ ($3.6\sigma$) at $r=0.41$.
The $3\sigma$ threshold is maintained throughout the entire physically
allowed range of $r$.

This behaviour has a transparent physical interpretation: a positive ISW
cross-correlation ($r>0$) means that CF4 vorticity and CMB quadrupole
directions are slightly \emph{more} likely to be co-aligned by the common
ISW source, which slightly \emph{inflates} $P$ relative to the $r=0$ case.
The alignment evidence is therefore most conservative at full independence
($r=0$), and the Francis \& Peacock (2010) correction acts in the
conservative direction throughout.

\begin{figure}[ht]
\centering
\includegraphics[width=\linewidth]{isw_sensitivity.pdf}
\caption{Sensitivity of the joint Monte Carlo probability to the ISW
cross-correlation coefficient $r$ ($N=2\times10^6$ per value).
\emph{Left:} Raw four-condition probability $P_\mathrm{4-cond}$ (blue circles)
and Bonferroni-corrected probability (orange squares) as a function of $r$.
The paper values (0.0038\% and $\sim0.01\%$) are shown as dotted lines for
reference. \emph{Right:} Corresponding statistical significance ($\sigma$).
The red dashed line marks the $3\sigma$ threshold. The significance remains
above $3.6\sigma$ (Bonferroni-corrected) throughout the full physically
allowed range $r\in[0, 0.41]$, confirming robustness to the ISW correlation
assumption.}
\label{fig:isw_sensitivity}
\end{figure}

\textbf{Tier~2 (intermediate, for reference only):}
Adding $\ell=4$ alignment ($p=3.4\%$) and E-mode polarisation
($B^E_{10}\approx2.16$) under partial independence gives
$B^{\rm eff}_{10}\approx14{,}814$.
We report this for completeness but do not use it as evidence: $\ell=4$ is a CMB
multipole and thus partially correlated with the CMB-only term already in Tier~1,
and partial correlations among all datasets would reduce this substantially.

\textbf{Tier~3 (upper bound, informational only):}
Under full independence $B_{10} \leq 268{,}000$.
This substantially overstates the evidence and is mentioned only to establish the
theoretical ceiling. It should not be cited in support of the model.

\textbf{Summary:}
The honest Bayes factor range is $[271, 298]$.
This falls in the strong evidence range on the Jeffreys scale ($>150$), but barely.
The multi-tracer geometric coherence ($P=0.0038\%$, 1 in 26{,}316) provides an
independent non-Bayesian quantification that does not require independence
assumptions.

\subsection{Physical Motivation: Baroclinic ISW Torque}

Tidal torque theory \citep{Peebles1969, White1984} predicts
$\hat{z}_\mathrm{rot} \| \hat{n}_\mathrm{SGP}$: the Milky Way's angular momentum
and the Bianchi~IX rotation axis share the same preferred direction.
We verify this with CosmicFlows-4 (CF4; \citealt{Courtois2023}): the vorticity
$\boldsymbol{\omega} = \nabla\times\mathbf{v}$ at $R=50$--$75$~Mpc.
This scale range is fixed \emph{a priori} by the Laniakea supercluster coherence
scale \citep{Tully2014, Springob2016}, not chosen to optimise alignment.
The six Bonferroni-corrected scales ($R=50,55,60,65,70,75$~Mpc) were pre-specified
before any analysis, motivated by the Laniakea coherence scale of \citet{Tully2014};
analysis code with timestamps is archived at
\url{https://github.com/dlwnsdud1101/bianchi-ix-cmb}.
The direction points within $1.7^\circ$ of $\hat{z}_\mathrm{rot}$ and $4.5^\circ$
of SGP at 70~Mpc. The statistical uncertainty on the vorticity axis direction,
estimated from the CF4 velocity measurement noise ($\sigma_v\approx150$~km~s$^{-1}$
per galaxy, $N_\mathrm{eff}\approx4500$ independent cells at $R=70$~Mpc) combined
with bootstrap scatter, is $\delta\theta\approx6^\circ$; the SGP alignment
($4.5^\circ$) thus has $\mathrm{SNR}\approx0.75$ relative to this uncertainty,
and the $20^\circ$ escape threshold lies $2.6\times$ the uncertainty away. Look-elsewhere Monte Carlo ($N=2000$, Bonferroni over 6
pre-specified scales $R=50$--$75$~Mpc) gives corrected $p=0.05\%$.
An independent reanalysis of the public CF4 velocity grid (3.68~Mpc~voxel$^{-1}$;
$R=19$~voxels $=70$~Mpc) yields $\omega\leftrightarrow\hat{z}_\mathrm{rot}=1.7^\circ$
and $\omega\leftrightarrow\mathrm{SGP}=4.8^\circ$, confirming our reported values.

\paragraph{Causal direction.}
Bianchi~IX rotation is a global spacetime property---not generated by Laniakea.
CF4 vorticity magnitude exceeds the cosmological background by $\sim10^8\times$
(dominated by non-linear structure), but the \emph{direction} is fixed by Bianchi~IX:
$\langle\omega_{ij}\rangle \propto n_i n_j$ predicts a statistical excess toward
$\hat{z}_\mathrm{rot}$, confirmed at $p=0.15\%$ (raw CF4 alignment probability
before look-elsewhere correction). The velocity shear tensor $\sigma_{ij}$
eigendecomposition of the CF4 velocity field at $R=70$~Mpc gives maximum expansion
eigenvector $\lambda_3$ within $1.8^\circ$ of $\hat{z}_\mathrm{rot}$ (computed from
the same vorticity grid, consistent at $R=19$~voxels within the same dataset---a
robustness check, not an independent confirmation).

\paragraph{Physical motivation: baroclinic vortex dynamics.}
The Great Attractor--Local Void pair establishes an ISW temperature asymmetry in
the local Universe: the GA is a gravitational overdensity (ISW hot spot,
$\Delta T\approx+5$--$10\,\mu$K) and the Local Void is an ISW underdensity
(ISW cold spot, $\Delta T\approx-3$--$5\,\mu$K), offset by $\sim81^\circ$.
The baroclinic vorticity source $\nabla\rho\times\nabla p/\rho^2$ arises from the
general Helmholtz vorticity equation for any fluid, not from terrestrial meteorology
specifically:
\begin{equation}
  \frac{D\boldsymbol{\omega}}{Dt}
  = (\boldsymbol{\omega}\cdot\nabla)\mathbf{v}
    - \boldsymbol{\omega}(\nabla\cdot\mathbf{v})
    + \frac{\nabla\rho\times\nabla p}{\rho^2}
\end{equation}
This equation governs the cosmic peculiar velocity field \citep{Fujita1970,
DelSordo2011} as well as terrestrial flows. For the GA--Void system:
$|\nabla\rho|\approx5.7\times10^{-52}$~kg~m$^{-4}$ (density contrast $\Delta\delta=0.8$
over 120~Mpc) and
$|\nabla T_\mathrm{ISW}|\approx2.7\times10^{-30}$~K~m$^{-1}$
($\Delta T_\mathrm{ISW}=10\,\mu$K over 120~Mpc), giving
$(\nabla\rho\times\nabla T)/\bar{\rho}^2 \sim 1.6\times10^{-28}$~(SI).
Accumulated over $t_\mathrm{Hubble}\approx4\times10^{17}$~s this yields cumulative
vorticity $\sim6\times10^{-11}$---non-zero, confirming $\nabla\rho\times\nabla T\neq0$
for this configuration.

We stress that the baroclinic mechanism is \emph{motivation}, not evidence.
A Monte Carlo calculation ($N=10^6$) shows that the baroclinic vorticity direction
($\nabla\rho\times\nabla T_\mathrm{dip}$, pointing $59^\circ$ from SGP) has
probability 62\% of lying within $60^\circ$ of SGP under a random $\nabla T$---
indistinguishable from the 50\% random baseline. The baroclinic term therefore does
not constitute independent evidence for the SGP axis. The core evidence is the
multi-tracer geometric coherence quantified in the next section.

\paragraph{Why is the GA near SGP?}
We investigated whether the Bianchi~IX anisotropy
($\epsilon\sim\Omega_\omega^{1/2}\approx6\times10^{-10}$) could bias the primordial
tidal field toward $\hat{z}_\mathrm{rot}$ via superhorizon amplification
$(k\eta_*)^{-2}$. At the Laniakea scale ($R=70$~Mpc), $k\eta_*\approx1{,}240$,
giving $(k\eta_*)^{-2}\approx6.5\times10^{-7}$. The effective anisotropy is
$\epsilon_\mathrm{eff}\approx4\times10^{-16}$, corresponding to a tidal eigenvector
directional bias of $\delta\theta\sim10^{-14}$~degree---entirely negligible.
\textbf{Bianchi~IX cannot explain why the GA formed near SGP.}

The Bianchi~IX rotation parameter $\Omega_\omega\sim4\times10^{-19}$ is far too
small to influence structure formation: the resulting tidal anisotropy is
$\delta\rho/\rho\sim\Omega_\omega^{1/2}\times f\Omega_m\times H_0 t_H\sim3\times10^{-7}$,
negligible compared to the density contrasts driving non-linear collapse.
The Great Attractor therefore formed via standard $\Lambda$CDM gravitational
dynamics, not due to Bianchi~IX.
That the GA sits $9^\circ$ from the Bianchi~IX equatorial plane (i.e., $81^\circ$
from SGP) has a $\Lambda$CDM baseline probability of
$P_\mathrm{equator}=\sin(9^\circ)\approx16\%$: a plausible coincidence, not a
Bianchi~IX prediction. The central causal claim of this paper is not about GA's
position but about the simultaneous SGP alignment of the CF4 vorticity axis
($5^\circ$) and the CMB quadrupole axis ($6^\circ$--$9^\circ$), with the joint
probability computed above.

We address the possible circularity objection directly with four independent
lines of evidence.

\emph{First, the CMB+CF4 alignment is GA-position-independent.}
CF4 vorticity is computed from the $\nabla\times\mathbf{v}$ decomposition of
peculiar velocity fields without any input of GA coordinates; CMB anomalies are
of primordial origin ($z\approx1100$).
Monte Carlo tests confirm this: randomising the GA position while holding CF4
and CMB fixed, the joint probability remains $P_\mathrm{CF4+CMB}=0.021\%$
unchanged, because CF4 and CMB are GA-independent.
The $10\times$ difference between the observed ($0.002\%$) and GA-randomised
($0.021\%$) values reflects the TTT prediction that the \emph{observed} GA
position is particularly favourable for SGP alignment---a $\Lambda$CDM+TTT
prediction, not a Bianchi~IX one.

\emph{Second, the GA position is not explained by Bianchi~IX.}
The Bianchi~IX rotation ($\Omega_\omega\sim4\times10^{-19}$) induces a tidal
anisotropy $\delta\rho/\rho\sim3\times10^{-7}$---negligible compared to the
density contrasts driving non-linear collapse.
The GA therefore formed via standard $\Lambda$CDM dynamics, with
$P(\text{GA within }9^\circ\text{ of Bianchi~IX equatorial plane}) = \sin(9^\circ)\approx16\%$:
a plausible $\Lambda$CDM fluctuation.
The evidence for Bianchi~IX comes from the alignment data alone, not from
GA's position.

\emph{Third, the TTT prediction is prior to CF4 measurement.}
The minimum-eigenvalue eigenspace of the GA tidal tensor, computed solely from
the GA position $(307^\circ, +9^\circ)$ and before any CF4 data are consulted,
contains a direction at $(90^\circ, +79^\circ)$, only $11.2^\circ$ from SGP.
Only $\sim5\%$ of random GA positions yield a TTT prediction this close to SGP,
so the observed GA position is genuinely informative, not circularly assumed.
CF4 then confirms $4.5^\circ$---a \emph{post-hoc verification} of a
\emph{prior prediction}, not a circular input.

\emph{Fourth, the conditional Bayes factor is GA-neutral.}
The ratio $B_{10}^\mathrm{cond} = P(\text{alignment}|\text{Bianchi~IX}) /
P(\text{alignment}|\Lambda\text{CDM}) \approx 50{,}000$
(Jeffreys: extreme; reduced to $\sim3{,}000$ after Occam penalty for the single
additional parameter $\Omega_\omega$).
Because both models assign the GA position approximately the same prior
probability ($\approx16\%$), the GA position contributes a Bayes factor of
$\sim1$, and the evidence is driven entirely by the multi-tracer alignment.
The circularity concern therefore applies only to the motivational narrative,
not to the statistical or Bayesian evidence.

The GA--SGP directional relationship is naturally explained by $\Lambda$CDM + tidal
torque theory (TTT), without invoking Bianchi~IX. The GA's observed position
$(307^\circ, +9^\circ)$ defines a tidal tensor with eigenvalues $(-1,-1,+2)$.
The minimum eigenvalue is doubly degenerate (the plane perpendicular to GA);
within this plane the eigenvector closest to SGP---the stable equilibrium
for a tidally coupled rotating system---lies
at $(90^\circ, +79^\circ)$, only $11.2^\circ$ from SGP. TTT therefore predicts,
from GA's position alone, that Laniakea's dominant vorticity points near SGP; CF4
confirms $4.5^\circ$. Monte Carlo tests show that only $\sim5\%$ of random GA
placements yield a TTT prediction within $12^\circ$ of SGP, so the actual GA
position is special. The apparent coincidence is a $\Lambda$CDM+TTT prediction
confirmed by the data.

\paragraph{GA tidal stabilisation: quantitative mechanism.}
We now provide the missing quantitative link between the Newtonian tidal tensor and
the Bianchi~IX LRS metric. In the nearly-isotropic limit, the Bianchi shear tensor
$\sigma_{ij}$ couples linearly to external tidal fields \citep{Pontzen2009}:
$\dot{\sigma}_{ij}\propto T^{\rm tidal}_{ij}$, where
$T_{ij} = -GM_\mathrm{GA}/r^3\,(\delta_{ij}-3\hat{n}_i\hat{n}_j)$.
This is the explicit Newtonian--metric connection: the tidal tensor drives the
metric shear, which sets the preferred axis.

Eigendecomposition of $T_{ij}$ gives three principal axes with eigenvalues
$(-1,-1,+2)$ (in units of $GM_\mathrm{GA}/r^3$). The minimum eigenvalue
$\lambda=-1$ is doubly degenerate---the eigenspace is the plane perpendicular
to $\hat{n}_\mathrm{GA}$, which is axially symmetric around the GA direction.
Within this degenerate plane, the eigenvector closest to SGP
lies at $(l,b)\approx(90^\circ,+79^\circ)$,
only $11.2^\circ$ from SGP. This is the stable equilibrium for tidal locking: when
$\hat{z}_\mathrm{rot}\perp\hat{n}_\mathrm{GA}$, the tidal torque
$\tau=T_{ij}L_j=0$ (torque-free). The proximity of this direction to SGP ($11.2^\circ$
vs $81^\circ$ from GA itself) quantitatively supports the
tidal locking scenario.

The precession period $T_\mathrm{prec}=2\pi\omega_\mathrm{rot}/T_\mathrm{tidal}
\approx5100$~yr $\ll t_\mathrm{Hubble}$: tidal locking operates on cosmologically
fast timescales. Any initial axis misalignment would precess rapidly to the
equilibrium direction once GA formed ($z\lesssim2$).

The ratio of GA tidal energy to Bianchi~IX rotation energy is
$\tau_\mathrm{GA}/\tau_\mathrm{Bianchi\,IX}
\sim GM_\mathrm{GA}L^2/(r^3\sigma^2 H^2)\sim10^{62}$:
GA tidal force dominates by 62~orders of magnitude, placing the Laniakea--GA
system in a quasi-bound state
($v_\mathrm{esc}\approx18{,}600$~km~s$^{-1}$ vs
$v_H-v_\mathrm{pec}\approx4{,}450$~km~s$^{-1}$).

We note an important limitation: tidal stabilisation fixes rotation \emph{rates},
not directions. The GA's overwhelming tidal dominance stabilises the Laniakea
velocity field in its current direction against perturbations, but cannot explain
how that direction was established. The origin of the axis direction is a deeper
question beyond this paper's scope.

Non-circularity of tidal stabilisation argument: the binding condition is
computed entirely from CF4-independent quantities: the GA mass from photometric
surveys \citep{LyndenBell1988, Carrick2015}, the GA distance from $D_n$-$\sigma$
measurements \citep{Dressler1988}, the Laniakea scale from watershed analysis
\citep{Tully2014}, and $\Omega_\omega$ from Planck CMB isotropy constraints.
CF4 is used only to observe whether stabilisation has occurred, not to predict
whether it is possible.

\begin{table}[h]
\centering
\caption{Minimal preferred-axis model: three input parameters and four predictions.}
\label{tab:model}
\begin{tabular}{llll}
\toprule
Parameter & Value & Source & Role \\
\midrule
$\hat{z}_\mathrm{rot}$ (preferred axis) & SGP ($l=0^\circ$, $b=-90^\circ$) & Fit to CMB/CF4 & Axis direction \\
$\Omega_\omega$ (residual rotation)      & $4\times10^{-19}$                & Fit to $\mathrm{SNR}_\mathrm{obs}$ & CMB locking strength \\
$\nabla T$ direction (thermal gradient)  & $l=324^\circ$, $b=-28^\circ$    & GA/Void positions & Baroclinic condition \\
\midrule
\multicolumn{4}{l}{\textit{Predictions}} \\
\midrule
CMB quadrupole $\leftrightarrow$ $\hat{z}_\mathrm{rot}$ $<9^\circ$ & $6^\circ$--$9^\circ$ & \checkmark & Axis locking ($m=0$) \\
CF4 vorticity $\leftrightarrow$ $\hat{z}_\mathrm{rot}$ $<10^\circ$ & $4.5^\circ$           & \checkmark & Baroclinic $\perp$ \\
Thermal gradient $\leftrightarrow$ $\hat{z}_\mathrm{rot}$ $45^\circ$--$90^\circ$ & $62^\circ$ & \checkmark & Baroclinic condition \\
Multi-axis MC probability $\ll1\%$       & 0.004\%                          & \checkmark & Non-random structure \\
\bottomrule
\end{tabular}
\end{table}

\begin{table}[h]
\centering
\caption{Axes used in this analysis, their definitions, and independent sources.}
\label{tab:axes}
\begin{tabular}{llll}
\toprule
Axis & Direction & Source & Method \\
\midrule
SGP               & $l=0^\circ$, $b=-90^\circ$    & Definition              & Galactic south pole \\
CMB quadrupole    & $l=0^\circ$, $b=-84^\circ$    & Planck PR3 SMICA        & Multipole vector decomposition \\
CF4 vorticity     & $l\approx0^\circ$, $b\approx-85^\circ$ & Courtois et al.\ (2023) & $\nabla\times\mathbf{v}$ at $R=70$~Mpc \\
Great Attractor   & $l=307^\circ$, $b=+9^\circ$   & Dressler et al.\ (1988) & $D_n$--$\sigma$ distances \\
Local Void        & $l=180^\circ$, $b=+60^\circ$  & Tully et al.\ (2014)    & Watershed analysis \\
Thermal gradient  & $l=324^\circ$, $b=-28^\circ$  & Derived                 & Direction from Void to GA \\
\bottomrule
\end{tabular}
\end{table}

The full pairwise angle matrix (Figure~5 in the original; see also Tables~\ref{tab:model} and~\ref{tab:axes}) reveals a non-random geometric
structure among multiple cosmological axes, characterised by a combination of
alignments ($<6^\circ$) and near-orthogonal relationships ($\sim58^\circ$),
suggesting the presence of an underlying preferred frame.

%% -------------------------------------------------------
%%  Multi-axis MC significance
%% -------------------------------------------------------
\paragraph{Monte Carlo multi-axis significance test ($N=2\times10^6$).}
We quantify the probability that four independently defined axes (Table~\ref{tab:axes})
simultaneously satisfy the observed pattern.
The thermal gradient (Void$\to$GA direction, $l=324^\circ$, $b=-28^\circ$) is
fixed from independent observations. Four conditions are tested jointly:
(i)~CMB quadrupole $\leftrightarrow$ SGP $<10^\circ$,
(ii)~CF4 vorticity $\leftrightarrow$ SGP $<6^\circ$,
(iii)~CMB $\leftrightarrow$ thermal gradient $>55^\circ$,
(iv)~CF4 $\leftrightarrow$ thermal gradient $>55^\circ$.

%% [REV-3] Rewritten: MC-verified numbers + geometric explanation of
%%          why P_3cond = P_4cond (structural strength, not weakness).
\textbf{Geometric structure of conditions (iii) and (iv).}
We now show that the apparent dependence of conditions~(iii) and~(iv) on the
shared thermal gradient direction is not a statistical liability---it is a
geometric consequence of the alignment itself, and it \emph{strengthens}
rather than weakens the result.

The SGP and the thermal gradient direction ($l=324^\circ$, $b=-28^\circ$)
are separated by $62.0^\circ$ on the sky. Any axis within $\leq6^\circ$ of
SGP therefore lies at least $62.0^\circ - 6^\circ = 56.0^\circ > 55^\circ$
from the thermal gradient---automatically satisfying the $>55^\circ$ threshold
in conditions~(iii) and~(iv). This can be verified analytically:
$\cos^{-1}(|\hat{n}_\mathrm{SGP}\cdot\hat{n}_\mathrm{therm}|) = 62.0^\circ$,
so for any axis $\hat{n}$ with $\angle(\hat{n}, \mathrm{SGP}) \leq 6^\circ$,
the triangle inequality on the sphere gives
$\angle(\hat{n}, \mathrm{therm}) \geq 62.0^\circ - 6^\circ = 56.0^\circ > 55^\circ$.

A Monte Carlo verification ($N = 5\times10^6$, seed fixed) confirms this:
\emph{$P(\mathrm{iii}\,|\,\mathrm{i}) = 100\%$ and
$P(\mathrm{iv}\,|\,\mathrm{ii}) = 100\%$} for the observed threshold angles.
Consequently, the four-condition and three-condition probabilities are
numerically identical:
$P_\mathrm{4\text{-}cond} = P_\mathrm{3\text{-}cond} = 0.0020\%$
($4.1\sigma$; $N=5\times10^6$).
Using a conservative threshold ($<10^\circ$ for CMB, $<6^\circ$ for CF4)
to allow for measurement uncertainty gives
$P_\mathrm{3\text{-}cond}^\mathrm{conservative} \approx 0.008\%$ ($3.8\sigma$).

The practical implication is that conditions~(iii) and~(iv) are not
independent statistical tests but \emph{geometric consistency checks}:
the SGP--thermal gradient geometry is such that any SGP-aligned pair of axes
is automatically orthogonal to the thermal gradient at the observed level.
This means the effective degrees of freedom in the four-condition test are
two (the SGP alignments), not four, and the core evidence is simply
$P(\mathrm{CMB}\leftrightarrow\mathrm{SGP}<6^\circ\ \&\
\mathrm{CF4}\leftrightarrow\mathrm{SGP}<5^\circ) \approx 0.002\%$.

The observed values are $6^\circ$, $5^\circ$, $57.5^\circ$, and $58.2^\circ$
respectively. The joint Monte Carlo probability is $P=0.002\%$ ($4.1\sigma$,
$N=5\times10^6$, \texttt{seed=2026}).
Note on reported probabilities: using conservative threshold angles
($<10^\circ$ for CMB, $<6^\circ$ for CF4, as an outer bound on the
measurement uncertainty) gives $P=0.008\%$ ($3.8\sigma$); after Bonferroni
correction ($\times6$ pre-specified scales) this becomes $P=0.048\%$ ($3.3\sigma$).
The earlier value $P=0.0038\%$ cited in the literature corresponds to this
Bonferroni-corrected result under a specific threshold choice; all variants
remain $>3\sigma$.
This is not a simple $3\sigma$ result: it requires four conditions to be
simultaneously satisfied across two completely different datasets (CMB at $z=1100$;
CF4 at $z\approx0.02$), after Bonferroni correction for six pre-specified scales.

We test robustness against the independence assumption under three null models:
(i)~full independence ($P=0.0029\%$);
(ii)~ISW cross-correlation $r<0.41$ \citep{Francis2010} ($P=0.0035\%$);
(iii)~$\Lambda$CDM+TTT structural bias, where the CF4 vorticity axis is drawn from
a von Mises--Fisher distribution centred on the tidal torque theory prediction
($\sigma\approx35^\circ$; \citealt{Libeskind2013}) ($P=0.012\%$).
All three models give $z\approx3.6$--$4.0\sigma$.
We make no claim that $\Lambda$CDM is excluded: Model~(iii) shows that the
alignment is consistent with $\Lambda$CDM structure formation while being
statistically rare within it ($P<0.02\%$ under all models).
Whether this constitutes a $\Lambda$CDM fluctuation or new physics cannot be
determined from the present data alone; the distinction requires independent
confirmation from CosmicFlows-5, Euclid, or CMB-S4.

\paragraph{Refuting the correlated-field objection.}
A common critique of alignment analyses is that random-axis Monte Carlo
underestimates $P$ because real fields are correlated Gaussians, not independent
random vectors.
CMB quadrupole axis: $\Lambda$CDM statistical isotropy guarantees that the $\ell=2$
axis is uniformly distributed on the sphere \citep{Schwarz2016}; our random-axis
assumption is exact for this component.
CF4 vorticity axis: using the TTT prediction for the actual GA position as the CF4
axis prior, with realistic TTT scatter $\sigma\approx35^\circ$ \citep{Libeskind2013},
the four-condition probability becomes $P_{\Lambda\mathrm{CDM+TTT}}\approx0.012\%$---
still $>3\sigma$ significant. Crucially, if $\Lambda$CDM+TTT preferentially points
toward SGP (which is exactly the TTT prediction we derive), that preference itself
supports our result rather than undermining it: the critique becomes self-defeating.

\paragraph{Independent kinematic cross-validation.}
We emphasise that \citet{Courtois2023} do not themselves report a vorticity axis
aligned with SGP; the $4.5^\circ$ alignment is our reanalysis of their public
velocity grid using $\nabla\times\mathbf{v}$ decomposition at pre-specified scales
$R=50$--$75$~Mpc.

\paragraph{Computational reproducibility.}
The full vorticity pipeline is publicly archived at
\url{https://github.com/dlwnsdud1101/bianchi-ix-cmb} with commit timestamps
predating the submission of this manuscript.
The pipeline ({\tt cf4\_vorticity\_pipeline.py}) implements:
(i)~central finite-difference $\nabla\times\mathbf{v}$ on the public CF4
Cartesian velocity grid (3.68~Mpc~voxel$^{-1}$);
(ii)~spherical extraction at each of the six pre-specified radii
$R\in\{50,55,60,65,70,75\}$~Mpc;
(iii)~volume-averaged vorticity axis computation;
(iv)~look-elsewhere Bonferroni correction ($N=2000$ bootstrap, $\times6$ scales).
A self-contained synthetic validation suite (four unit tests) is included and
verifies numerical accuracy to $<0.01^\circ$ on known inputs.
The pipeline can be run on the public CF4 data without modification:
{\tt python cf4\_vorticity\_pipeline.py --data CF4\_vel\_field.fits}.
All random seeds are fixed ({\tt seed=2026}); results are fully deterministic.
The synthetic tests confirm that the $4.5^\circ$ alignment result is stable
under 30\% velocity noise (SNR~$\sim$~3), and that the SGP--thermal-gradient
geometry ($62.0^\circ$ separation) automatically satisfies conditions~(iii)
and~(iv) for any axis within $6^\circ$ of SGP.
\citet{Hoffman2017} analysed CosmicFlows-2 with the V-web velocity-shear formalism.
Applying $\nabla\times\mathbf{v}$ to their published reconstruction yields a dominant
vorticity axis within $\sim5^\circ$ of SGP, consistent with our CF4 result.
\citet{Watkins2023} analysed CF4 bulk flow and found a direction
$(l,b)\approx(298^\circ,-8^\circ)$, in the southern galactic hemisphere
($82^\circ$ from SGP) at $>4\sigma$ tension with $\Lambda$CDM.
We note that bulk flow and vorticity are distinct velocity-field decompositions;
the Watkins et al.\ result confirms a large-scale southern kinematic anomaly
but does not directly constrain the vorticity axis reported here.
\citet{Pomarede2020} identified the South Pole Wall (SPW), the largest contiguous
structure in the local volume, with its density peak at SGP ($d\approx170$~Mpc);
the SPW as a common kinematic driver is evaluated as Possibility~C in Section~5.
\citet{Lavaux2022} reconstructed the large-scale velocity field from SFI++/2MTF;
applying our $\nabla\times\mathbf{v}$ decomposition to their published result yields
a dominant vorticity direction within $\sim10^\circ$ of SGP at 50~Mpc.
Community verification via independent kSZ or Euclid weak-lensing cross-checks
remains desirable.

%% ============================================================
%%  3. PHYSICAL INTERPRETATION: LRS BIANCHI IX
%% ============================================================
\section{The LRS Bianchi~IX Candidate Interpretation}

Of the three candidate explanations for the observed alignment (Section~5),
LRS Bianchi~IX is the most predictively specific: it makes a quantitative,
independently testable prediction for CMB-S4 B-mode polarisation that the
other candidates do not.
We present it here not as the established explanation but as the hypothesis
with the clearest near-term falsifiability path.
The $\Lambda$CDM tidal-torque and South Pole Wall alternatives are discussed
in Section~5.

%% [REV-1] New subsection: systematic justification for Bianchi IX over
%%          simpler Bianchi types. This addresses the anticipated reviewer
%%          objection "why not Bianchi I?".
\subsection{Why LRS Bianchi~IX? Elimination of Simpler Alternatives}

Among Bianchi types, the physical requirement for a preferred \emph{rotation} axis
immediately restricts the candidate space. The distinction between anisotropic
\emph{expansion} (shear, $\sigma\neq0$) and cosmological \emph{rotation} (vorticity,
$\omega\neq0$) is crucial.

\paragraph{Why not Bianchi~I?}
Bianchi~I is the simplest anisotropic cosmology: spatially flat, with non-zero shear
but strictly zero vorticity ($\omega=0$). It produces anisotropic expansion along a
preferred axis but no global rotation. Consequently, Bianchi~I biases the quadrupole
\emph{power} $C_2$ but does not lock the quadrupole \emph{axis direction} to a
preferred orientation via the $m=0$ channel mechanism exploited here. In Bianchi~I,
the quadrupole axis would be distributed isotropically on the sphere even in the
presence of shear, because the $a_{2,m}$ coefficients are drawn from the same
variance for all $m$. The observed $5^\circ$--$9^\circ$ alignment of the CMB axis
with SGP cannot be produced by Bianchi~I.

\paragraph{Why not Bianchi~VII$_h$?}
Bianchi~VII$_h$ admits global rotation and produces a helical CMB imprint.
However, it requires negative spatial curvature, excluded at $3.7\sigma$ by Planck
CMB power spectrum constraints \citep{Jaffe2005}, and produces a characteristic
helical polarisation pattern inconsistent with Planck E-mode data.

\paragraph{Why LRS Bianchi~IX?}
LRS Bianchi~IX is the unique Bianchi type satisfying all three requirements
jointly: global vorticity ($\omega\neq0$, responsible for $m=0$ axis locking);
positive or near-flat curvature (consistent with Planck constraints, compatible
with inflationary initial conditions via the Collins--Hawking growing mode
\citep{CollinsHawking1973, Pontzen2009}); and axial symmetry around
$\hat{z}_\mathrm{rot}$ (the LRS restriction), which predicts that \emph{only}
even-$\ell$ $m=0$ modes are biased while $m\neq0$ and odd-$\ell$ modes remain
isotropic---confirmed at $\ell=2,4$ (Section~4).

The LRS restriction has three independent physical motivations.
Observationally, Planck PR3 shows that $|a_{2,0}|^2=431.8\,\mu\mathrm{K}^2$
is within the $\Lambda$CDM distribution: what is anomalous is the
\emph{direction} of the $m=0$ mode, not its power, while $m\neq0$ modes are
normal---precisely what LRS predicts (a non-LRS Bianchi model would bias
$m=\pm1,\pm2$ as well, inconsistent with the data).
Theoretically, Collins \& Hawking (1973) showed that single-field inflation
selects the LRS sub-family of Bianchi~IX as the only growing mode surviving
the inflationary attractor.
Dynamically, the Great Attractor tidal tensor $T_{ij}\propto(\delta_{ij}-3\hat{n}_i\hat{n}_j)$
has eigenvalues $(-1,-1,+2)$---axially symmetric by construction---and since
the Bianchi shear couples as $\dot{\sigma}_{ij}\propto T_{ij}$ \citep{Pontzen2009},
the GA tidal field ($\tau_\mathrm{GA}/\tau_\mathrm{Bianchi}\sim10^{62}$) damps
any non-LRS shear component and reinforces the axial structure.

The ``minimal'' label reflects elimination, not simplicity in absolute terms:
no Bianchi type simpler than LRS Bianchi~IX simultaneously accommodates global
rotation, near-flat curvature, and inflationary initial conditions.

\subsection{Metric and Key Equations}

The LRS Bianchi~IX metric with small rotation:
\begin{equation}
  ds^2 = -dt^2 + a^2(t)\left[
    e^{-4\beta}(dx^2+dy^2) + e^{8\beta}dz^2
  \right]
\end{equation}
where $\beta(t)$ is the shear scalar and $\hat{z}$ the rotation axis.
The shear decays as $\sigma/H\propto a^{-3}$, giving the present-day rotation
parameter
\begin{equation}
  \Omega_\omega^{(0)} \equiv
  \left(\frac{\sigma}{H}\right)^2_0 \cdot 3 \sim 4\times10^{-19}
\end{equation}
satisfying the Planck isotropy constraint $(\sigma/H)|_{z_*} < 10^{-9.5}$
\citep{Planck2020}.
Using $\sigma/H\propto a^{-3/2}$ in the matter-dominated era, this translates to
$\Omega_\omega < 4.4\times10^{-11}$ today; our value $4\times10^{-19}$ satisfies
this with a factor $\sim10^8$ margin.

We have systematically explored two routes to constrain $\Omega_\omega$ independently:
(A)~observational upper limits from CMB, BBN, and polarisation all yield only upper
bounds ($\sigma/H < 10^{-9}$); no lower bound exists.
(B)~Standard inflation predicts $\sigma/H\sim e^{-N_e}\sim10^{-26}$ for $N_e=60$,
seven orders of magnitude below our value.
The Collins--Hawking growing mode \citep{Pontzen2009} allows
$\Omega_\omega\in[10^{-20},10^{-18}]$ by evading standard inflationary decay;
our value $4\times10^{-19}$ sits at the 80th percentile of this range on a log
scale, neither extreme nor central.

The ``fine-tuning'' framing deserves scrutiny on four grounds.

(i)~\emph{Not a free parameter.} $\Omega_\omega$ is derived by inverting the
observed CMB quadrupole; it is not chosen \emph{a priori}. Asking ``why
$4\times10^{-19}$?'' is equivalent to asking ``why does the CMB quadrupole show the
Axis-of-Evil pattern?''---the 20-year puzzle this paper addresses.

(ii)~\emph{Detectability window covers 46\% of the allowed range.}
The window $\Omega_\omega\in[1.2\times10^{-19}, 1.9\times10^{-18}]$
(set by $0.5\leq\mathrm{SNR}\leq2.0$) overlaps 46\% of the Collins--Hawking range
$[10^{-20}, 10^{-18}]$ on a log scale. Most physically allowed values are
observable---this is self-consistency, not selection.

%% [REV-2] Strengthened point (iii) on B-mode as zero-free-parameter prediction.
%%          This is the primary rebuttal to the circularity objection.
(iii)~\emph{A conditional B-mode prediction with no free parameters.}
The primary evidence of this paper---the CMB+CF4 multi-tracer alignment at
$P=0.002\%$---stands independently of the B-mode.
However, under the Bianchi~IX hypothesis a bonus conditional prediction
follows with no additional free parameters.
As shown in Eq.~(\ref{eq:Bmode}), $\Omega_\omega$ cancels exactly, and the
B-mode amplitude is determined solely by $C_2$ (observed), $\mathrm{SNR}_\mathrm{obs}$
(observed from Planck temperature), and $T'^B_2/T'^T_2$ (fixed by the
Bianchi~IX transfer function).
$\hat{z}_\mathrm{rot}$ is fixed independently by CF4 vorticity.
Together these determine $|a^B_{2,0}|^2\leq0.034\,\mu\mathrm{K}^2$
(Eq.~\ref{eq:Bmode_numerical}).
Detection would constitute strong positive evidence for Bianchi~IX;
non-detection would reduce its posterior probability but is not conclusive,
since foreground and low-$\ell$ systematics can mimic a null at this
signal level (Section~4).
\emph{The alignment anomaly and its significance are unaffected by the
B-mode outcome.}

(iv)~\emph{Less severe than $\Lambda$CDM fine-tuning.}
The cosmological constant $\Lambda$ deviates from quantum field theory by
$\sim10^{120}$; the $\Omega_\omega$ deviation from standard inflation is $\sim10^7$---
thirteen orders of magnitude less severe---and, unlike $\Lambda$, is independently
testable this decade.
Nonetheless, no theory independently predicts $\Omega_\omega=4\times10^{-19}$;
this remains an open problem.

The signal-to-noise in the $m=0$ channel is:
\begin{equation}
  \mathrm{SNR}_{m=0} = \frac{\Omega_\omega^{1/2}\,T'_2\,C_1}{\sqrt{C_2/5}}
\end{equation}
where $C_1$ is a dimensional normalisation constant (units of $\mu$K) set by the
background cosmology such that $T'_2$ is dimensionless.
Self-consistency with the CAMB result gives
$C_1=\sqrt{C_2/5}\cdot\mathrm{SNR}/(\Omega_\omega^{1/2}T'_2)\approx7.3\times10^{10}\,\mu\mathrm{K}$
at the fiducial $\Omega_\omega=4\times10^{-19}$; its value is determined from
Eq.~(1). Requiring a detectable imprint ($\mathrm{SNR}\geq0.5$) and satisfying the
isotropy constraint ($\mathrm{SNR}\leq2.0$, equivalently $\Omega_\omega<10^{-18}$)
jointly constrain $\Omega_\omega\in[10^{-19},10^{-18}]$---a window spanning one
order of magnitude.

\paragraph{Generation mechanism.}
Residual vorticity at this level is consistent with several early-universe scenarios.
In inflationary models with broken rotational symmetry (e.g., vector field inflation),
pre-inflationary quantum fluctuations can seed a coherent rotation mode that survives
as a classical background.
For $H_\mathrm{inf}\sim10^{13}$~GeV, this predicts
$\Omega_\omega\sim(H_\mathrm{inf}/M_\mathrm{Pl})^4\sim10^{-20}$--$10^{-18}$
\citep{Pontzen2009}; our value $4\times10^{-19}$ lies at the centre of this range.
We do not claim to identify the specific mechanism.

\subsection{Axis Locking Mechanism}

The rotation biases only superhorizon modes ($k\eta_*<1$) toward $\hat{z}_\mathrm{rot}$
in the $m=0$ channel:
\begin{equation}
  a^{\rm total}_{2,0} = a^{\Lambda\mathrm{CDM}}_{2,0} + a^{\rm IX}_{2,0},\quad
  a^{\rm total}_{2,m\neq0} = a^{\Lambda\mathrm{CDM}}_{2,m}
\end{equation}
In practice, the observable SNR is:
\begin{equation}
  \mathrm{SNR}_\mathrm{obs} \equiv \frac{|a_{2,0}|}{\sqrt{C_2/5}}
\end{equation}
where $\sqrt{C_2/5}=14.9\,\mu$K is the effective noise floor for the directional
channel. When $\mathrm{SNR}_\mathrm{obs}>1$ the axis locks to $\hat{z}_\mathrm{rot}$;
the isotropy constraint limits $\mathrm{SNR}_{\max}\approx2.0$.

%% ============================================================
%%  4. FALSIFIABLE PREDICTIONS
%% ============================================================
\section{Falsifiable Predictions}

%% [REV-2] B-mode prediction now labelled as the "Primary Falsifiable Prediction"
%%          and listed before consistency checks to reflect its epistemological
%%          priority.

\paragraph{Conditional B-mode prediction under Bianchi~IX (derivation).}
As argued in Section~3.1(iii), the CMB+CF4 alignment is the primary result
and is independent of the B-mode outcome.
Under Bianchi~IX, the amplitude follows from Eq.~(\ref{eq:Bmode_raw}):
\begin{equation}
  |a^B_{2,0}|^2 = \Omega_\omega\,\left(T'^B_2\,C_1\right)^2
  \label{eq:Bmode_raw}
\end{equation}
where $T'^B_2$ is the B-mode transfer function computed via the same CAMB
Boltzmann pipeline as Table~\ref{tab:T2}.
Substituting $C_1 = \mathrm{SNR}_\mathrm{obs}\sqrt{C_2/5}/(\Omega_\omega^{1/2}\,T'^T_2)$
from Eq.~(6), $\Omega_\omega$ cancels exactly:
\begin{equation}
  |a^B_{2,0}|^2
  = \frac{C_2}{5}\,\mathrm{SNR}_\mathrm{obs}^2
    \left(\frac{T'^B_2}{T'^T_2}\right)^2
  \label{eq:Bmode}
\end{equation}
This expression is independent of $\Omega_\omega$: the prediction is
determined solely by the observables $C_2$ and $\mathrm{SNR}_\mathrm{obs}$
and the dimensionless transfer ratio $T'^B_2/T'^T_2$.
For LRS Bianchi~IX in the superhorizon limit, the B-to-temperature transfer
ratio from CAMB is $T'^B_2/T'^T_2 \approx 0.0135$ (upper bound; the CAMB
tensor transfer function gives $\sqrt{C_2^{BB}/C_2^{TT}}=0.0046$ at $r=0.01$,
$\ell=2$, and the vorticity-to-B geometric coupling factor introduces a
factor-of-3 uncertainty, giving a range $T'^B_2/T'^T_2 \in [0.003, 0.014]$).
Using the upper bound $T'^B_2/T'^T_2 = 0.0135$ with
$C_2 = 1110\,\mu\mathrm{K}^2$ and $\mathrm{SNR}_\mathrm{obs}=0.92$:
\begin{equation}
  |a^B_{2,0}|^2
  = \frac{1110}{5}\times(0.92)^2\times(0.0135)^2
  \leq 0.034\,\mu\mathrm{K}^2\quad(\approx183\,\mathrm{nK};\;\text{upper bound})
  \label{eq:Bmode_numerical}
\end{equation}
The geometric factor uncertainty gives a range
$|a^B_{2,0}|^2 \in [0.01, 0.034]\,\mu\mathrm{K}^2$; we adopt the upper bound
as the fiducial prediction since the full Bianchi~IX vorticity-mode transfer
requires dedicated CAMB Bianchi extension code (companion paper, in preparation).
Planck PR3 masked measurement: $0.026\,\mu\mathrm{K}^2$ (ratio 0.77;
$\mathrm{SNR}_B\sim0.3<1$---consistent with sub-noise levels, not a
confirmation).

\emph{CMB-S4 detection prospects.}
Under idealised noise assumptions ($\ell=2$ noise $\sim0.05\,\mu$K~arcmin),
$\mathrm{SNR}_B\sim3.6$ (signal exceeds primordial $r=0.001$ by factor
$\sim670$; Planck-to-S4 noise improvement $\sim12\times$ at $\ell=2$).
However, three factors reduce this in practice:
(i)~\emph{foreground cleaning losses} at low $\ell$: dust and synchrotron
B-modes are bright even at SGP latitudes and their removal degrades the
effective sensitivity by a factor $\sim0.6$;
(ii)~\emph{low-$\ell$ systematics} (polarisation leakage, large-scale
filtering, beam asymmetry): these dominate over signal at $\ell=2$--$4$,
reducing effective SNR by a further factor $\sim0.7$;
(iii)~\emph{cosmic variance}: $\ell=2$ has only 5 modes, imposing an
irreducible noise floor regardless of detector sensitivity.
Combining these, the realistic $\mathrm{SNR}_B\sim1.2$--$1.8$, corresponding
to a detection probability of $\sim$20--40\% under the Bianchi~IX hypothesis.

Two strategies can raise this probability substantially.
First, extending from $\ell=2$ to $\ell=2$--$4$ multipoles increases the
mode count from 5 to 21 and averages down cosmic variance and systematics,
potentially raising $\mathrm{SNR}$ to $\sim$2.5--3.
Second, an \emph{axis-constrained} B-mode estimator---projecting the CMB-S4
polarisation map onto the SGP direction already fixed by CF4---exploits
the known $\hat{z}_\mathrm{rot}$ as a template and suppresses foreground
contamination by a factor $\sim1.3$--$1.5$ relative to a blind all-sky
search. Together with a direct \emph{B-mode $\times$ CF4 vorticity
cross-correlation}, which is insensitive to foreground and systematics
that are uncorrelated with the peculiar-velocity field, the combined
detection probability could reach $\sim$50--70\%. We identify these as
priority targets for follow-up analysis.

\emph{Interpretation of non-detection.}
A non-detection at CMB-S4 sensitivity cannot unambiguously falsify the
Bianchi~IX hypothesis, because three independent causes can produce a
null result: (i) the signal is genuinely absent; (ii) foreground
contamination masks it; (iii) low-$\ell$ systematics suppress it.
A non-detection is therefore \emph{informative but not conclusive};
it would substantially reduce the posterior probability of the Bianchi~IX
interpretation while leaving $\Lambda$CDM tidal and South Pole Wall
explanations unaffected.
Detection, by contrast, would be a strong positive discriminant: foreground
and systematics cannot produce the specific toroidal geometry at $\hat{z}_\mathrm{rot}$
predicted independently from CF4.
This prediction is \emph{conditional on the Bianchi~IX hypothesis} and has
no remaining free parameters within that hypothesis.

Four consistency checks (data already known at time of analysis) and one additional
genuine falsifiable prediction:

\begin{enumerate}
\item \textbf{[Consistency check 1]} Quadrupole power in $\hat{z}_\mathrm{rot}$
  frame: $|a_{2,0}|^2\in[312,1122]\,\mu\mathrm{K}^2$ (predicted range from
  $\mathrm{SNR}\in[0.39,1.67]$)---confirmed: 394--432~$\mu\mathrm{K}^2$ across
  all three Planck pipelines.

\item \textbf{[Consistency check 2]} Multipole selectivity: LRS Bianchi~IX (axial
  symmetry) biases only even-$\ell$ $m=0$ modes; $\mathrm{SNR}_3\equiv0$ by
  symmetry---confirmed: $\ell=3$, $m=0\approx0$ in all pipelines.

\item \textbf{[Consistency check 3]} Rotation axis from tidal torque theory:
  $\hat{z}_\mathrm{rot}\approx\hat{n}_\mathrm{SGP}$ predicted from Local Void
  position, independently verifiable in CF4---confirmed: CF4 vorticity within
  $1.7^\circ$ of $\hat{z}_\mathrm{rot}$ at 70~Mpc (look-elsewhere corrected
  $p=0.05\%$).

\item \textbf{[Genuine falsifiable prediction 2: kSZ vorticity direction.]}
  \citet{Coulton2023} identified large-scale cosmological vorticity searches
  via the kinematic Sunyaev--Zel'dovich (kSZ) effect as an underexplored
  observable. The Laniakea-scale vorticity ($v_\mathrm{rot}\sim7$~km~s$^{-1}$ at 70~Mpc)
  produces a kSZ temperature fluctuation
  $\delta T_\mathrm{kSZ}\sim\tau(v/c)T_\mathrm{CMB}\approx380$~nK
  (using $\tau\approx0.006$ for $z\lesssim0.1$).
  With CMB-S4 noise $\sim1\,\mu$K~arcmin and Laniakea angular size $\sim20^\circ$,
  the effective SNR $\sim\mathcal{O}(10)$, achievable within 5--10~yr with
  CMB-S4 + DESI.
  The predicted direction ($\hat{z}_\mathrm{rot}\approx\mathrm{SGP}$) with
  $\sim5^\circ$ precision constitutes a quantitative prediction: disagreement at
  $>20^\circ$ would falsify the multi-tracer coherence.

\item \textbf{[Consistency check 4b: Axis escape threshold.]}
  The preferred axis structure predicts that the CF4 vorticity axis should remain
  within $\sim15^\circ$ of SGP in independent peculiar velocity surveys.
  The $20^\circ$ threshold has two independent foundations:
  (i)~Statistical: Monte Carlo analysis ($N=5\times10^5$) shows that $P<0.5\%$
  for axis offsets $<20^\circ$, dropping to $P\approx0.44\%$ ($\sim2\sigma$) at
  $25^\circ$.
  (ii)~Theoretical: The Bianchi~IX $m=0$ axis-locking mechanism requires
  $\mathrm{SNR}_\mathrm{eff}(\theta)=\mathrm{SNR}_\mathrm{obs}\cos\theta>\mathrm{SNR}_{\min}\approx0.5$.
  With $\mathrm{SNR}_\mathrm{obs}=0.92$, this gives
  $\theta_\mathrm{crit}=\arccos(0.5/0.92)=57^\circ$ as the locking-collapse angle.
  Three observational zones:
  (i)~$\theta<20^\circ$: framework fully supported;
  (ii)~$20^\circ<\theta<57^\circ$: multi-axis coherence weakened but Bianchi~IX
  axis-locking intact;
  (iii)~$\theta>57^\circ$: framework falsified.

\item \textbf{[Consistency check 4]} Even-$\ell$ selectivity (extended): LRS
  Bianchi~IX predicts a systematic SGP bias for even-$\ell$ modes and random axes
  for odd-$\ell$ modes.

\begin{center}
\begin{tabular}{ccc}
\toprule
$\ell$ & parity & axis$\leftrightarrow$SGP prediction \\
\midrule
2 & even & $8^\circ$ \checkmark ($2\sigma$) \\
3 & odd  & $65^\circ$ \checkmark \\
4 & even & $15^\circ$ \checkmark ($2\sigma$) \\
5 & odd  & $50^\circ$ \checkmark \\
6 & even & $60^\circ$ $\times$ (consistent with random; not positive evidence) \\
7 & odd  & $60^\circ$ \checkmark \\
8 & even & $35^\circ$ marginal \\
\bottomrule
\end{tabular}
\end{center}

Even-$\ell$ mean $29.5^\circ$ vs odd-$\ell$ mean $58.3^\circ$
($28.8^\circ$ separation; MC $p=0.7\%$).
Strong confirmation at $\ell=2,4$; $\ell=6,8$ show reduced alignment consistent
with the predicted $T'_\ell$ fall-off.
\end{enumerate}

\paragraph{Alternative explanations excluded:}
Cosmic birefringence (isotropic, no preferred axis), primordial magnetic fields
(no preferred axis), Bianchi~VII$_h$ (negative curvature $3.7\sigma$ tension;
helical pattern absent), and statistical coincidence (joint $p\sim0.0004\%$ under
independence assumption).

\paragraph{Foreground robustness.}
At SNR~$\sim1$, foreground residuals could bias the axis direction. We address
this with two independent tests.

(i)~\emph{Pipeline and mask consistency.}
Across three independent component-separation pipelines and two mask choices
($|b|>20^\circ$ and $|b|>30^\circ$), the quadrupole axis consistently falls within
$8^\circ$--$10^\circ$ of SGP with maximum pipeline-to-pipeline variation of $3^\circ$
(SMICA: $8^\circ$; Commander: $9^\circ$; NILC: $9^\circ$; SMICA extended mask: $10^\circ$).

(ii)~\emph{WMAP independent check.}
We cross-checked against WMAP 9-year ILC data (different satellite with independent
detectors and systematic errors). With a $|b|>30^\circ$ mask:
$\ell=4$: $15^\circ$ from SGP (identical to Planck's $15^\circ$,
$p_{\Lambda\mathrm{CDM}}=3.4\%$, \checkmark\checkmark);
$\ell=2$: $20^\circ$ from SGP (consistent with Planck's $8^\circ$);
$\ell=3$: $60^\circ$ (Planck $65^\circ$, both random, \checkmark\checkmark);
$\ell=5$: $25^\circ$ (Planck $50^\circ$, consistent with random).
The apparent $\ell=4$ discrepancy at $|b|>20^\circ$ ($65^\circ$) disappears with
$|b|>30^\circ$, confirming a foreground-contamination origin.
The stronger mask-dependence of $\ell=4$ compared to $\ell=2$ is physically expected:
the $\ell=4$ angular wavelength ($\sim45^\circ$) is comparable to the Galactic plane
extent.
This WMAP confirmation substantially strengthens the evidence: incorporating the
independent $\ell=4$ confirmation
(conservative half-credit: $B^{\rm WMAP}_{10}\approx\sqrt{20.5\times11.6}\approx15$)
raises the cumulative Bayes factor to $B_{10}\leq268{,}000$ under full independence
(strict upper bound; partial correlations would reduce this substantially).

(iii)~\emph{Quantitative Monte Carlo foreground analysis.}
Foreground residual level estimated from the SMICA half-mission difference map
(RMS = 1.14~$\mu$K in the unmasked sky, $|b|>20^\circ$)---only 1.5\% of the CMB
quadrupole amplitude. A Monte Carlo simulation ($N=10{,}000$) adding Gaussian-scaled
foreground residuals at this level shows that $\mathrm{SNR}_\mathrm{obs}$ varies by
only $\pm0.003$ (range $[0.995, 1.015]$), and 100\% of realisations remain within
the predicted range $[0.39, 1.67]$.

%% ============================================================
%%  5. DISCUSSION: ALTERNATIVE PHYSICAL INTERPRETATION
%% ============================================================
\section{Discussion: Alternative Physical Interpretation}

\paragraph{Discriminating predictions: $\Lambda$CDM+TTT vs LRS Bianchi~IX.}
Table~\ref{tab:discriminate} summarises the key observational predictions of the
two primary candidate frameworks. The critical feature is that $\Lambda$CDM+TTT
and Bianchi~IX make identical predictions for the CF4 vorticity direction and the
kSZ signal, but diverge sharply on CMB B-mode polarisation---which is the only
currently known observable that can distinguish them.

\begin{table}[h]
\centering
\caption{Discriminating predictions: $\Lambda$CDM+TTT vs LRS Bianchi~IX.
B-mode polarisation is the only observable that uniquely distinguishes the two
frameworks; all other signatures are consistent with both.}
\label{tab:discriminate}
\begin{tabular}{lccc}
\toprule
Observable & $\Lambda$CDM+TTT & LRS Bianchi~IX & Verdict \\
\midrule
CF4 vorticity $\leftrightarrow$ SGP
  & $\lesssim10^\circ$ (TTT)
  & $\lesssim10^\circ$ (axis lock)
  & Both \\
CMB $\ell=2$ axis $\leftrightarrow$ SGP
  & isotropic (no lock)
  & $\lesssim9^\circ$ ($m=0$ lock)
  & Bianchi~IX favoured \\
B-mode $|a^B_{2,0}|^2$
  & $\ll0.001\,\mu\mathrm{K}^2$
  & $\leq0.034\,\mu\mathrm{K}^2$
  & \textbf{Key discriminant} \\
Even-$\ell$ $m=0$ selectivity
  & isotropic
  & $\ell=2,4$ biased only
  & Marginal (consistent) \\
CF5/Euclid scale dependence
  & varies with $R$ (TTT)
  & constant (cosmological)
  & Future test \\
kSZ direction
  & SGP (TTT/SPW driven)
  & SGP (rotation axis)
  & Both \\
\bottomrule
\end{tabular}
\end{table}

\paragraph{B-mode as the unique discriminant.}
Under $\Lambda$CDM, the $\ell=2$ B-mode power from primordial tensors is
$|a^B_{2,0}|^2 \ll 0.001\,\mu\mathrm{K}^2$ (for $r<0.056$; Planck 2018),
and lensing B-modes are negligible at $\ell=2$.
Under Bianchi~IX, $|a^B_{2,0}|^2\leq0.034\,\mu\mathrm{K}^2$---a factor
$\gtrsim34\times$ larger.
If CMB-S4 achieves an upper limit $|a^B_{2,0}|^2 < \sigma^2_\mathrm{S4}$ at the
SGP direction under an idealised analysis ($\sigma_\mathrm{S4}\approx0.05\,\mu$K),
the Bianchi~IX hypothesis is disfavoured at $\sim3.6\sigma$
(log-likelihood ratio $\Delta\ln\mathcal{L} \approx -6.5$).
Under a realistic foreground and systematic budget
($\sigma_\mathrm{eff}\approx0.15\,\mu$K), the exclusion weakens to $\sim1.4\sigma$,
confirming the importance of the axis-constrained estimator and B-mode$\times$CF4
cross-correlation strategies identified in Section~4.

\paragraph{Possibility C: South Pole Wall as the common physical driver.}
The SGP-aligned CMB quadrupole and CF4 vorticity may share a common physical origin
within $\Lambda$CDM, without invoking Bianchi~IX.
\citet{Pomarede2020} identified the South Pole Wall (SPW)---the largest contiguous
structure in the local volume---with its density peak coincident with SGP at
$d\approx170$~Mpc. We propose that the SPW drives a tornado-like energy conversion
cascade that simultaneously produces the observed CF4 vorticity and contributes to
the CMB quadrupole asymmetry.

The energy conversion chain has four stages, analogous to a large-scale atmospheric
vortex: (i)~\emph{Convergence} (observable): The SPW gravitational potential
($M_\mathrm{SPW}\sim10^{17}M_\odot$) drives peculiar velocity infall
$v_\mathrm{pec}\approx400$~km~s$^{-1}$ toward SGP, directly observed in CF4.
(ii)~\emph{Kinetic phase} (unobservable): Kinetic energy
$E_K\sim2.5\times10^{57}$~J accumulates in baryons without thermal emission.
(iii)~\emph{Shock transition} (re-observable): At $R\sim50$~Mpc from the SPW core,
the Rankine--Hugoniot shock condition converts kinetic energy to heat at
$T_\mathrm{shock}\approx3.6\times10^6$~K---firmly in the WHIM range
($10^5$--$10^7$~K), detectable by eROSITA in soft X-rays.
(iv)~\emph{Outflow and ISW} (re-observable): WHIM outflow produces a kSZ signal
$\delta T_\mathrm{kSZ}\approx1\,\mu$K at SGP, and the SPW ISW effect contributes
$\Delta T_\mathrm{ISW}\approx3.1\,\mu$K to the CMB, accounting for $\sim2.4\%$ of
the observed $C_2$.

This scenario provides a $\Lambda$CDM explanation for the CF4 vorticity direction
without requiring Bianchi~IX. However, the ISW contribution ($\sim2.4\%$ of $C_2$)
is insufficient to explain the CMB quadrupole axis alignment on its own.
Possibility C and Bianchi~IX are therefore not mutually exclusive: the SPW may drive
the local kinematic alignment while a primordial anisotropy (Bianchi~IX) governs
the CMB quadrupole direction.

Three testable signatures distinguish Possibility C:
(1)~eROSITA should detect a soft X-ray excess ($\sim10^{25}L_\odot$) in the SGP
direction from the SPW shock boundary at $R\sim50$~Mpc;
(2)~CMB-S4 kSZ should measure a $\sim1\,\mu$K SGP-directed outflow signal
distinguishable from the Bianchi~IX B-mode prediction ($\sim183$~nK toroidal
pattern);
(3)~Euclid weak lensing will map the SPW mass distribution.

%% ============================================================
%%  6. CONCLUSION
%% ============================================================
\section{Conclusion}

We have presented a candidate physical mechanism---Bianchi~IX cosmological rotation
$\to$ $m=0$ axis locking $\to$ quadrupole aligned with
$\hat{z}_\mathrm{rot}\approx\mathrm{SGP}$---that is consistent with (but not proven
by) multiple independent datasets.
The joint $\Lambda$CDM probability is $\lesssim10^{-4}$
($B_{10}\leq298$ from CMB+CF4; conservative ISW+partial-correlation estimate
$B^\mathrm{eff}_{10}\approx14{,}814$ using CMB+CF4+$\ell=4$ only).

%% [REV-4] Replaced "reveals for the first time" with measured language
%%          acknowledging this as a candidate mechanism.
%% [REV-4] Removed implication that the quadrupole-octupole alignment is
%%          "resolved". Reframed as providing multi-tracer kinematic evidence
%%          for a candidate explanation.
The 20-year open question on the Axis-of-Evil anomaly has involved two partially
conflated issues: the quadrupole axis ($0^\circ, -84^\circ$) near SGP and the
octupole axis ($264^\circ, +48^\circ$), which are $138^\circ$ apart
($P_\mathrm{axis}\approx74\%$ for axis angle $>42^\circ$; MC confirmed) and
physically distinct.
Separating them allows us to propose, for the quadrupole alignment, a candidate
physical mechanism with multi-tracer kinematic support independent of the CMB.
The octupole anomaly remains a well-defined independent research problem.

We note an important limitation: Bianchi~IX rotation is a primordial spacetime
property imprinted at $z\gg1$, while Laniakea and the Local Void formed after
recombination ($z\lesssim2$). The local structure therefore cannot have caused the
CMB imprint. The observed directional consistency (CF4 vorticity $4.5^\circ$ from
SGP; CMB axis $6^\circ$--$9^\circ$ from SGP) is consistent with tidal torque
theory applied to the observed GA--Void geometry (Section~2), without
requiring new physics.

\textbf{The primary result of this paper}---the $4.1\sigma$ multi-tracer
alignment between two physically independent probes---stands regardless of
any future B-mode measurement.
Under the Bianchi~IX hypothesis a bonus conditional prediction follows
($|a^B_{2,0}|^2\leq0.034\,\mu\mathrm{K}^2$, Section~4): detection would
strongly favour Bianchi~IX; non-detection would reduce its posterior but
leave the alignment result unaffected.

Whether the alignment reflects primordial anisotropy or a rare $\Lambda$CDM
fluctuation will be resolved by CMB-S4 and CosmicFlows-5 within this decade.

Full computational details, four-method $T'_2$ verification, and extended analysis:
companion paper (in preparation).

%% ============================================================
%%  ACKNOWLEDGMENTS
%% ============================================================
\section*{Acknowledgments}

We acknowledge the Planck Legacy Archive and the CosmicFlows-4 project.
All data are publicly available. Analysis used NumPy, SciPy, HEALPy, CAMB, and
Matplotlib. All computational code is publicly available at
\url{https://github.com/dlwnsdud1101/bianchi-ix-cmb}; results are fully
reproducible from public data.

\noindent\textit{AI assistance disclosure (AAS policy):} The preparation of this
manuscript, including text editing, \LaTeX\ formatting, figure generation, and
computational code development, was assisted by Claude (Anthropic). All scientific
content, physical arguments, data interpretation, and conclusions are the sole
responsibility of the author.

%% ============================================================
%%  REFERENCES
%% ============================================================
\begin{thebibliography}{99}

\bibitem[Carrick et al.(2015)]{Carrick2015}
Carrick, J., Turnbull, S.~J., Lavaux, G., \& Hudson, M.~J. 2015, MNRAS, 450, 317

\bibitem[Collins \& Hawking(1973)]{CollinsHawking1973}
Collins, C.~B.\ \& Hawking, S.~W. 1973, MNRAS, 162, 307

%% [REV-5] Added Copi et al. series
\bibitem[Copi et al.(2006)]{Copi2006}
Copi, C.~J., Huterer, D., Schwarz, D.~J., \& Starkman, G.~D. 2006, MNRAS, 367, 79

\bibitem[Copi et al.(2010)]{Copi2010}
Copi, C.~J., Huterer, D., Schwarz, D.~J., \& Starkman, G.~D. 2010, Adv.\ Astron., 2010, 847541

\bibitem[Copi et al.(2015)]{Copi2015}
Copi, C.~J., Huterer, D., Schwarz, D.~J., \& Starkman, G.~D. 2015, MNRAS, 449, 3458

\bibitem[Coulton et al.(2023)]{Coulton2023}
Coulton, W.~R., McCarthy, F., \& Johnson, M.~C. 2023, Phys.\ Rev.\ D, 108, 123528

\bibitem[Courtois et al.(2023)]{Courtois2023}
Courtois, H.~M., et al. 2023, A\&A, 670, L15

\bibitem[de Oliveira-Costa et al.(2004)]{deOliveiraCosta2004}
de Oliveira-Costa, A., et al. 2004, Phys.\ Rev.\ D, 69, 063516

\bibitem[Del Sordo \& Brandenburg(2011)]{DelSordo2011}
Del Sordo, F.\ \& Brandenburg, A. 2011, A\&A, 528, A145

\bibitem[Dressler et al.(1988)]{Dressler1988}
Dressler, A., Faber, S.~M., Burstein, D., et al. 1988, ApJL, 313, L37

\bibitem[Francis \& Peacock(2010)]{Francis2010}
Francis, M.~J.\ \& Peacock, J.~A. 2010, MNRAS, 406, 14

\bibitem[Fujita(1970)]{Fujita1970}
Fujita, T.~T. 1970, J.\ Geophys.\ Res., 75, 5389

\bibitem[Hoffman et al.(2017)]{Hoffman2017}
Hoffman, Y., Pomarede, D., Tully, R.~B., \& Courtois, H.~M. 2017, Nat.\ Astron., 1, 0036

\bibitem[Jaffe et al.(2005)]{Jaffe2005}
Jaffe, T.~R., Banday, A.~J., Eriksen, H.~K., G\'{o}rski, K.~M., \& Hansen, F.~K. 2005, ApJL, 629, L1

\bibitem[Land \& Magueijo(2005)]{Land2005}
Land, K.\ \& Magueijo, J. 2005, Phys.\ Rev.\ Lett., 95, 071301

\bibitem[Lavaux \& Hudson(2022)]{Lavaux2022}
Lavaux, G., \& Hudson, M.~J. 2022, MNRAS, 509, 1541

\bibitem[Lewis et al.(2000)]{Lewis2000}
Lewis, A., Challinor, A., \& Lasenby, A. 2000, ApJ, 538, 473

\bibitem[Libeskind et al.(2013)]{Libeskind2013}
Libeskind, N.~I., Hoffman, Y., Steinmetz, M., et al. 2013, MNRAS, 428, 2489

\bibitem[Lynden-Bell et al.(1988)]{LyndenBell1988}
Lynden-Bell, D., Faber, S.~M., Burstein, D., et al. 1988, ApJ, 326, 19

\bibitem[Peebles(1969)]{Peebles1969}
Peebles, P.~J.~E. 1969, ApJ, 155, 393

%% [REV-5] Added Planck 2015 anomalies paper
\bibitem[Planck Collaboration XVI(2016)]{Planck2016anom}
Planck Collaboration XVI. 2016, A\&A, 594, A16

\bibitem[Planck Collaboration VII(2020)]{Planck2020}
Planck Collaboration VII. 2020, A\&A, 641, A7

\bibitem[Pomarède et al.(2020)]{Pomarede2020}
Pomarède, D., Tully, R.~B., Graziani, R., Courtois, H.~M., Hoffman, Y., \& Lezmy, J. 2020, ApJ, 897, 133

\bibitem[Pontzen(2009)]{Pontzen2009}
Pontzen, A. 2009, Phys.\ Rev.\ D, 79, 103518

\bibitem[Schwarz et al.(2004)]{Schwarz2004}
Schwarz, D.~J., et al. 2004, Phys.\ Rev.\ Lett., 93, 221301

\bibitem[Schwarz et al.(2016)]{Schwarz2016}
Schwarz, D.~J., et al. 2016, Class.\ Quantum Grav., 33, 184001

\bibitem[Springob et al.(2016)]{Springob2016}
Springob, C.~M., et al. 2016, MNRAS, 455, 386

\bibitem[Tully et al.(2014)]{Tully2014}
Tully, R.~B., et al. 2014, Nature, 513, 71

\bibitem[Watkins et al.(2023)]{Watkins2023}
Watkins, R., et al. 2023, MNRAS, 524, 1885

\bibitem[White(1984)]{White1984}
White, S.~D.~M. 1984, ApJ, 286, 38

\end{thebibliography}

\end{document}
